% ===================================================================
% THE ONE-PAGE SELF-ASSESSMENT — a real worksheet, not an image.
% Replaces \includegraphics{...book1-one-page.png}.
% Set as a full page with \clearpage either side.
% ===================================================================
\clearpage
\thispagestyle{empty}
\begin{center}
  {\large\bfseries The one-page self-assessment}\\[0.3em]
  {\small Fill this in before you open the second book.}
\end{center}
\vspace{0.5em}

\newcommand{\qline}{\par\vspace{0.35em}\noindent\rule{\linewidth}{0.4pt}}
\newcommand{\qblock}[3]{%
  \noindent{\bfseries #1}\par
  {\footnotesize\itshape #2}\par
  \vspace{0.3em}
  \qline #3
  \vspace{0.42em}
}

\qblock{1. Which engine pays me?}
  {Global (sells to the world), state (funded by government), or domestic
   (sells to the 5.9 million here). Say which, and who hands you the money.}
  {\qline}

\qblock{2. Which wave is my demand riding?}
  {Ageing, the fertility collapse, the shrinking household, or the remaking of
   care. If none of them, say so plainly.}
  {\qline}

\qblock{3. What is my industry, and what does the data say?}
  {Name the industry, its size, whether it is growing, and how good the public
   data on it is. If the data is thin, write ``thin'' --- that is an answer.}
  {\qline}

\qblock{4. Am I facing a wall or a door?}
  {Concentrated (a few names hold most of it --- a wall) or fragmented
   (nobody owns it --- a door). Count the ones your customer could name.}
  {\qline}

\qblock{5. Which two or three words look open?}
  {From the brand map: which words do the incumbents already own, and which
   are sitting free? Name two or three that look genuinely open.}
  {\qline\qline}

\vspace{0.6em}
\noindent{\footnotesize\itshape
This page may be photocopied. If you cannot answer question five, that is what the second book is for.}
\clearpage
